\documentclass{article}

\usepackage{neurips_2026}

\usepackage[utf8]{inputenc}
\usepackage[T1]{fontenc}
\usepackage{hyperref}
\usepackage{url}
\usepackage{booktabs}
\usepackage{amsfonts}
\usepackage{nicefrac}
\usepackage{microtype}
\usepackage{xcolor}
\usepackage{amsmath,amssymb}
\usepackage{multirow}
\usepackage{graphicx}

\title{Same-Class Batch Normalization Statistics Are Worse Than Wrong-Class:\\A Forensic Study of an Asymmetric Substitution Effect}

\author{%
  Anonymous Author(s)\\
  Affiliation\\
  Address\\
  \texttt{email}
}

\begin{document}

\maketitle

\begin{abstract}
Replacing the global batch-normalization (BN) running statistics of a trained convolutional network with statistics estimated only from samples of one class is a benign-sounding intervention---class-conditional statistics are, after all, just the natural per-class refinement of the global ones. We show this intervention has a sharp and counter-intuitive structure. On CIFAR-10 with a SmallResNet, substituting class-$c$ statistics when evaluating class-$c$ samples reduces accuracy from 91.1\% to 3.5\%, while substituting a deliberately mismatched class's statistics leaves 65.4\% accuracy. The same-class condition is thus roughly $26\times$ more damaging than the wrong-class one, and a simple ``out-of-distribution input'' explanation cannot account for the asymmetry. The ordering replicates across four architectures (SmallResNet, VGG-11-BN, SimpleCNN, WideResNet-28-10) and three datasets (CIFAR-10, CIFAR-100, Tiny-ImageNet-200). Two additional observations rule out the most natural alternative explanations: (i) a linear probe on the post-substitution penultimate features achieves $\geq 99.7\%$ accuracy on CIFAR-10 (and 100\% on the larger-label-space datasets), so the representation itself is not damaged; (ii) GroupNorm models, which have no running statistics to swap, are unaffected; and (iii) temperature scaling reduces the resulting expected calibration error from 0.394 to 0.151 but recovers exactly 0\% of the lost accuracy, so the failure is not a logit-scale issue. A linear interpolation between class-conditional and global statistics yields a smooth monotone recovery in classifier accuracy but a non-monotone dip in probe accuracy near $\alpha\approx 0.6$. We treat these findings as a forensic observation about BN-statistic substitution, and discuss interpretations.
\end{abstract}

\section{Introduction}
\label{sec:intro}

Batch normalization \citep{ioffe2015batch} stores per-layer running mean and variance vectors that are used at inference to standardize activations. The running statistics are typically thought of as a benign approximation to the true population mean and variance: a numerical convenience that affects optimization but not what the network represents.

We test that picture by perturbing the running statistics in a controlled way. For a network trained on $C$ classes, we estimate one set of running statistics per class by passing the corresponding subset of the training data through the network in BN \texttt{train()} mode without updating any parameters. We then run inference under four conditions: \emph{global} (the standard buffers), \emph{same-class} (class-$c$ statistics for class-$c$ samples), \emph{wrong-class} (class-$(c+\lfloor C/2\rfloor) \bmod C$ statistics for class-$c$ samples), and \emph{random-class} (a random other class's statistics for each sample).

A skeptical reading of any single one of these conditions is straightforward: the classifier was trained against the global statistics, so swapping in any other statistics is a distribution shift, and accuracy drops are unsurprising. What this paper documents is a pattern that the distribution-shift reading does not predict. Same-class statistics are not a mild refinement of global statistics; they are the \emph{worst} of the three substitutions, and by a wide margin. On CIFAR-10 with a SmallResNet, same-class accuracy is 3.5\%---below the 10\% chance baseline---while wrong-class accuracy is 65.4\% and random-class accuracy is 59.0\%. The ordering ``same $\ll$ random $<$ wrong $\ll$ global'' replicates across four architectures and three datasets.

We document this asymmetry as the headline contribution and pair it with three controls that constrain its interpretation. (1) The penultimate features remain class-discriminative: a linear probe trained on the post-substitution features reaches $\geq 99.7\%$ accuracy on CIFAR-10 and 100\% on CIFAR-100 and Tiny-ImageNet-200, exceeding the global baseline. (2) GroupNorm \citep{wu2018group} variants of the same architecture, which carry no running statistics, are unaffected by construction. (3) Temperature scaling \citep{guo2017calibration} reduces the same-class expected calibration error from 0.394 to 0.151 but recovers exactly 0\% of the lost accuracy: the failure is in the logit \emph{ranking}, not its scale. A continuous interpolation between class-conditional and global statistics shows a smooth monotone recovery in classifier accuracy, but the linear probe is non-monotone with a minimum near $\alpha\approx 0.6$.

\paragraph{Contributions.}
\begin{enumerate}
\setlength{\itemsep}{1pt}
\item \textbf{An asymmetric substitution ordering.} Same-class BN statistics are roughly $26\times$ more damaging than wrong-class on CIFAR-10 (3.5\% vs 65.4\%) and the ordering same $\ll$ random/wrong $\ll$ global replicates across four architectures (SmallResNet, VGG-11-BN, SimpleCNN, WideResNet-28-10) and three datasets (CIFAR-10, CIFAR-100, Tiny-ImageNet-200).
\item \textbf{The features survive.} A linear probe on the post-substitution penultimate features achieves 99.9\% accuracy on CIFAR-10 (above the global baseline of 90.9\%), and 100\% on CIFAR-100 and Tiny-ImageNet-200. The substitution does not damage the representation.
\item \textbf{Two ruled-out alternative explanations.} The effect requires BN's running-statistic mechanism (GroupNorm models retain 81.6\% accuracy) and is not a logit-scale problem (temperature scaling halves ECE but recovers 0\% accuracy).
\item \textbf{A non-monotone interpolation.} Linearly interpolating from class-conditional to global statistics produces a smooth monotone recovery of classifier accuracy but a non-monotone probe-accuracy curve, indicating that the two endpoints correspond to qualitatively different geometric regimes for the activation space.
\end{enumerate}

We treat the paper as a forensic empirical study of a specific intervention, not a sweeping mechanistic claim about what BN statistics ``are''; Section~\ref{sec:discussion} discusses interpretations.

\section{Background and Related Work}
\label{sec:bg}

\paragraph{Batch normalization.} BN \citep{ioffe2015batch} normalizes activations using mini-batch statistics during training and exponentially-averaged \emph{running} statistics $(\mu_\ell,\sigma_\ell^2)$ at inference, after which a learned affine transform $\gamma_\ell, \beta_\ell$ is applied. \citet{santurkar2018does} and \citet{bjorck2018understanding} argue BN's main effect is to smooth the optimization landscape rather than to address ``internal covariate shift''. These works examine BN's training-time role; we examine the inference-time role of the running buffers specifically.

\paragraph{BN statistics as adaptable quantities.} A number of methods treat BN statistics as something that can be re-estimated rather than treated as fixed. AdaBN \citep{li2018revisiting} replaces them with target-domain statistics for domain adaptation; TENT \citep{wang2021tent} and \citet{schneider2020improving} update them on test data for robustness to distribution shift. \citet{frankle2020training} show that training only the BN affine parameters of a randomly-initialized network can reach surprisingly high accuracy. These works modify BN parameters in directions that \emph{help}; we instead document a regime in which a seemingly benign modification is sharply harmful.

\paragraph{Linear probing.} Linear probes \citep{alain2016understanding,kornblith2019similarity} measure whether class structure is linearly accessible in a representation. We use them in a counterfactual role: a network that mispredicts but whose features still admit a high-accuracy linear probe has a problem in the head, not in the encoder.

\paragraph{Calibration.} \citet{guo2017calibration} introduced expected calibration error (ECE) and temperature scaling. We use both as diagnostics. Temperature scaling is monotone in logits and so cannot change the argmax; if it nonetheless recovered accuracy, the same-class failure would be a mere over-confidence problem, which it is not.

\paragraph{Normalization variants.} GroupNorm \citep{wu2018group} computes statistics over feature-map groups within each example and stores no running averages; we use it as an architectural control. Normalization-free architectures \citep{brock2021high} replace BN entirely.

\section{Methodology}
\label{sec:method}

\paragraph{Architectures.} We train four image classifiers on CIFAR-10 with standard augmentation (random crop pad 4, horizontal flip), SGD with momentum 0.9, weight decay $5{\times}10^{-4}$, initial LR 0.1, cosine annealing. \textbf{SmallResNet}: a four-stage residual net with widths $[32,64,128,256]$ and 20 BN layers (10 seeds, 20 epochs); a five-stage \textbf{SmallResNet64} variant with widths $[32,64,128,256,512]$ is used for $64{\times}64$ inputs. \textbf{VGG-11-BN} (8 BN layers, 3 seeds), \textbf{SimpleCNN} (4-layer plain CNN, 4 BN layers, 3 seeds), \textbf{WideResNet-28-10} \citep{zagoruyko2016wide} (depth 28, widen factor 10, dropout $p{=}0.3$, 3 seeds). We additionally train SmallResNet on CIFAR-100 (5 seeds, 25 epochs) and SmallResNet64 on Tiny-ImageNet-200 (3 seeds, 25 epochs). Full hyperparameters are in Appendix~\ref{app:config}.

\paragraph{Class-conditional statistics.} Let $f_\theta$ be a trained CNN with BN layers $\{B_\ell\}$ carrying running means $\mu_\ell$ and running variances $\sigma_\ell^2$. For class $c$ we reset the BN buffers, place the model in \texttt{train()} mode (so BN updates the buffers via its standard EMA), iterate in \texttt{torch.no\_grad()} over the class-$c$ training subset, and snapshot the resulting buffers as $(\mu_\ell^{(c)},(\sigma_\ell^{(c)})^2)$. \emph{No model parameters are updated.} The original buffers are restored before the next class's pass.

\paragraph{Inference conditions.} For test sample $x$ with label $c$ we evaluate four conditions: \emph{global} (load original buffers), \emph{same-class} (load $(\mu_\ell^{(c)},(\sigma_\ell^{(c)})^2)$), \emph{wrong-class} (load class-$(c+\lfloor C/2\rfloor) \bmod C$), and \emph{random-class} (load a uniformly random class $\neq c$, drawn per-sample). For each condition we record top-1 accuracy, mean true-class softmax confidence, and the penultimate-layer features.

\paragraph{Linear probe.} Penultimate features (post global average pooling) are split 80/20; we fit ridge regression ($\lambda{=}10^{-3}$, one-hot targets, closed form) for CIFAR-10, and multinomial logistic regression (L-BFGS, $C{=}1.0$) for $\geq 100$-class settings, and report held-out accuracy.

\paragraph{Interpolation.} For $\alpha \in \{0.0,0.1,\dots,1.0\}$ we set $\mu_\alpha = \alpha \mu_{\text{global}} + (1-\alpha)\mu_{\text{same}}$ and $\sigma_\alpha^2 = \alpha\sigma^2_{\text{global}} + (1-\alpha)\sigma^2_{\text{same}}$ and evaluate.

\paragraph{Calibration controls.} ECE uses 15 equal-width bins. Temperature scaling \citep{guo2017calibration} optimizes a single scalar $T$ on a held-out validation set under NLL, then divides logits by $T$. Because $T>0$ preserves the argmax, any accuracy change under temperature scaling is by construction zero; the diagnostic value lies in comparing ECE before and after.

\paragraph{GroupNorm control.} We retrain SmallResNet with all BatchNorm layers replaced by GroupNorm (8 groups, fewer when channels are smaller). GroupNorm has no running statistics, so the substitution paradigm is inapplicable and the relevant question is whether the architecture suffers any analogous loss.

\paragraph{Statistical protocol.} SmallResNet/CIFAR-10 uses 10 seeds (42--51); SmallResNet/CIFAR-100 uses 5 (42--46); all other configurations use 3 (42--44). We report mean $\pm$ std and paired $t$-tests where stated.

\section{Results}
\label{sec:results}

\subsection{The asymmetric substitution ordering}
\label{sec:ordering}

\begin{table}[t]
\caption{SmallResNet on CIFAR-10 across the four BN-statistic conditions. Mean $\pm$ std over 10 seeds.}
\label{tab:primary}
\centering
\small
\begin{tabular}{@{}lccc@{}}
\toprule
Condition & Accuracy (\%) & Confidence (\%) & Linear Probe (\%) \\
\midrule
Global (baseline) & $91.08 \pm 0.12$ & $88.88 \pm 0.15$ & $90.86 \pm 0.55$ \\
Same-class        & $\mathbf{3.55 \pm 0.69}$  & $4.80 \pm 0.52$  & $99.88 \pm 0.09$ \\
Wrong-class       & $65.37 \pm 4.69$ & $50.64 \pm 3.22$ & $100.00 \pm 0.00$ \\
Random-class      & $58.99 \pm 5.84$ & $45.47 \pm 5.19$ & $97.63 \pm 1.11$ \\
\bottomrule
\end{tabular}
\end{table}

Table~\ref{tab:primary} shows the four-condition comparison on our most-replicated configuration. Three observations follow.

First, every non-global condition lies strictly below global accuracy, as a generic distribution-shift account would predict. Second, the ordering among the three substitution conditions is the opposite of what such an account would predict: same-class statistics, which by any natural metric are \emph{closest} to global (they differ from global only by within-class variation, not between-class variation), are the most damaging. Same-class accuracy (3.5\%) lies below the 10\% chance baseline; wrong- and random-class accuracies (65.4\% and 59.0\%) sit much higher. Same-vs-global accuracy is $t(9)=-374.1, p=3.5\times 10^{-20}$; same-vs-wrong is $t(9)=-39.1, p<10^{-10}$.

Third, per-class deltas are uniformly large and negative across all 10 CIFAR-10 classes (Appendix~\ref{app:perclass}), ranging from $-0.97$ for \emph{automobile} to $-0.72$ for \emph{cat}; the collapse is not driven by one or two pathological classes.

\paragraph{Cross-architecture replication.} Table~\ref{tab:cross_arch} shows the ordering replicates across all four architectures we tested. In every case, same-class accuracy is well below wrong-class accuracy, and the ratio is at least $4\times$ for the BN-heavy networks.

\begin{table}[t]
\caption{Cross-architecture results on CIFAR-10 (mean $\pm$ std). The ratio same/wrong is $<0.18$ in every architecture.}
\label{tab:cross_arch}
\centering
\small
\begin{tabular}{@{}llcccc@{}}
\toprule
Architecture & BN layers & Global & Same-class & Wrong-class & Same/Wrong \\
\midrule
SmallResNet ($n=10$)  & 20   & $91.08 \pm 0.12$ & $3.55 \pm 0.69$  & $65.37 \pm 4.69$ & $0.054$ \\
VGG-11-BN ($n=3$)     & 8    & $85.44 \pm 1.28$ & $10.88 \pm 1.91$ & $62.96 \pm 1.25$ & $0.173$ \\
SimpleCNN ($n=3$)     & 4    & $83.17 \pm 0.23$ & $26.33 \pm 0.48$ & $78.10 \pm 0.46$ & $0.337$ \\
WRN-28-10 ($n=3$)     & many & $90.61 \pm 0.51$ & $11.12 \pm 0.90$ & $64.06 \pm 3.14$ & $0.174$ \\
\bottomrule
\end{tabular}
\end{table}

The same-class accuracy correlates roughly with the number of BN layers: SmallResNet (20 layers) drops the most, SimpleCNN (4 layers) the least. WideResNet-28-10 (a deep BN-heavy network with dropout $p{=}0.3$) shows a same-class accuracy of 11.1\%, somewhat higher than SmallResNet's 3.5\% despite many more BN layers; the dropout regularization may be partially buffering the head against activation perturbation.

\paragraph{Cross-dataset replication.} Table~\ref{tab:cross_data} extends the experiment to CIFAR-100 and Tiny-ImageNet-200 with SmallResNet/SmallResNet64. The accuracy collapse is even sharper in larger label spaces: same-class accuracy is below the chance baseline ($1/C$) on both datasets. With $C=100$ or $C=200$, all three substitution conditions fall to floor and the same-vs-wrong gap is no longer measurable, but the substitution effect itself is more severe.

\begin{table}[t]
\caption{Cross-dataset results with SmallResNet (CIFAR-10/100) and SmallResNet64 (Tiny-ImageNet-200). Chance is $1/C$. Same-class accuracy lies \emph{below} chance on both larger datasets.}
\label{tab:cross_data}
\centering
\small
\begin{tabular}{@{}lcccccc@{}}
\toprule
Dataset & $C$ & Chance & Global & Same-class & Wrong-class & Probe (same) \\
\midrule
CIFAR-10           & 10  & 10.00\% & $91.08 \pm 0.12$ & $3.55 \pm 0.69$ & $65.37 \pm 4.69$ & $99.88 \pm 0.09$ \\
CIFAR-100          & 100 & 1.00\%  & $70.11 \pm 0.57$ & $0.60 \pm 0.49$ & $1.00 \pm 0.00$  & $100.00 \pm 0.00$ \\
Tiny-ImageNet-200  & 200 & 0.50\%  & $60.77 \pm 0.30$ & $0.17 \pm 0.24$ & $0.50 \pm 0.00$  & $100.00 \pm 0.00$ \\
\bottomrule
\end{tabular}
\end{table}

\subsection{The features survive: a probe-vs-classifier dissociation}
\label{sec:probe}

The right-most columns of Tables~\ref{tab:primary} and~\ref{tab:cross_data} report the linear probe accuracy under same-class statistics. Despite classification accuracy of 3.5\% on CIFAR-10 and 0.17\% on Tiny-ImageNet-200, a linear probe trained on the same penultimate features achieves 99.9\% and 100.0\% respectively---in the latter case, with zero variance across seeds.

This dissociation matters for interpreting Section~\ref{sec:ordering}: whatever the same-class substitution is doing, it is not destroying the encoder's class-discriminative information. A linear refit of the head on the same features is essentially perfect. The classifier's inability to recover the class is therefore a property of the \emph{specific learned head}, not of the representation it operates on. Notably the probe accuracy under same-class often \emph{exceeds} the global-statistic probe (99.9\% vs 90.9\% on CIFAR-10), suggesting that class-conditional normalization removes within-class variation in a way that makes the resulting features more linearly separable.

The same-class probe is somewhat lower for SimpleCNN (54.7\%; see Appendix~\ref{app:cross_arch}), the smallest architecture; on the BN-heavy networks (SmallResNet, VGG, WRN) the probe sits at $\geq 98.4\%$.

\subsection{Two ruled-out explanations}
\label{sec:controls}

\paragraph{Generic normalization-swap effect.} A natural alternative explanation is that the asymmetry is an artifact of \emph{any} deep architecture with running statistics that are perturbed at inference. We test this by training SmallResNet with all BatchNorm layers replaced by GroupNorm \citep{wu2018group}, which computes statistics from the current input and has no running buffers (3 seeds, 20 epochs). The GroupNorm model achieves $81.65\pm 1.94\%$ accuracy and a $81.08\pm 1.64\%$ linear probe under standard inference; there are no running statistics to substitute, so the asymmetry is inapplicable to it by construction. The effect is thus specific to the running-statistic mechanism in BN.

\paragraph{Logit-calibration issue.} A second alternative is that same-class statistics merely sharpen or rescale the logits in a way that softmax confidence and ECE pick up but argmax does not. We test this with temperature scaling \citep{guo2017calibration}: optimize a single scalar $T$ on validation NLL under same-class statistics, then divide all logits by $T$ at test time.

\begin{table}[t]
\caption{Temperature scaling on the same-class condition (SmallResNet, CIFAR-10, $n{=}3$ seeds). Temperature scaling is by construction monotone in logits and cannot change the argmax; the relevant comparison is the ECE before/after.}
\label{tab:tempscale}
\centering
\small
\begin{tabular}{@{}lcc@{}}
\toprule
Metric & Before ($T{=}1$) & After ($T \approx 3.43$) \\
\midrule
Accuracy (\%) & $2.87 \pm 0.49$ & $2.87 \pm 0.49$ \\
ECE           & $0.394 \pm 0.013$ & $0.151 \pm 0.003$ \\
\bottomrule
\end{tabular}
\end{table}

The optimum is $T \approx 3.43$: same-class logits are substantially over-scaled. Temperature scaling cuts ECE from 0.394 to 0.151---a substantial calibration improvement---but recovers 0\% of the lost accuracy. This is by construction (monotone-in-logits transforms preserve argmax), but the diagnostic content is the residual ECE: even after the optimal scalar rescaling the probabilities are still poorly calibrated. The under-the-hood logit \emph{ranking} has been disrupted, not just its scale. The same-class collapse is therefore not a temperature problem.

\subsection{A non-monotone interpolation between class-conditional and global}
\label{sec:interp}

We linearly interpolate the BN buffers between same-class ($\alpha=0$) and global ($\alpha=1$) and evaluate at each step. Table~\ref{tab:interp} shows four representative points; the full 11-point sweep is in Appendix~\ref{app:interp}.

\begin{table}[t]
\caption{Linear interpolation from same-class ($\alpha=0$) to global ($\alpha=1$) statistics on SmallResNet/CIFAR-10. Classifier accuracy rises monotonically; the linear probe is non-monotone, with a minimum near $\alpha\approx 0.6$. The full 11-point sweep is in Appendix~\ref{app:interp}.}
\label{tab:interp}
\centering
\small
\begin{tabular}{@{}cccl@{}}
\toprule
$\alpha$ & Accuracy (\%) & Probe (\%) & \\
\midrule
0.0 & 3.44  & 99.90 & class-conditional endpoint \\
0.3 & 23.14 & 94.85 & \\
0.6 & 65.17 & 74.65 & probe minimum \\
1.0 & 91.01 & 90.85 & global endpoint \\
\bottomrule
\end{tabular}
\end{table}

Two features are notable. First, classifier accuracy rises smoothly and monotonically from 3.4\% to 91.0\%---there is no phase transition. Second, the linear probe is \emph{non-monotone}: it is highest at the same-class endpoint (99.9\%), declines to a minimum of 74.65\% at $\alpha=0.6$, and partially recovers to 90.85\% at the global endpoint. Near $\alpha\approx 0.6$--$0.7$ the two curves cross.

The non-monotone probe trajectory is constraining for any explanation of the same-class effect. If the substitution simply moved features uniformly off-distribution, one might expect probe accuracy to be a monotone function of $\alpha$ (whether increasing or decreasing). Instead, the two endpoints both admit high-quality linear separability while an intermediate point does not, suggesting that the activation space at $\alpha=0$ and $\alpha=1$ correspond to two qualitatively different geometric regimes, with intermediate statistics partially mixing the two.

\section{Discussion}
\label{sec:discussion}

The four findings of Section~\ref{sec:results} together constrain how the same-class effect can be explained. Section~\ref{sec:controls} rules out two of the most natural alternatives: it is not a generic property of running-statistic perturbation (GroupNorm has no running statistics and is unaffected), and it is not a logit-scale issue (temperature scaling halves ECE without recovering accuracy). What remains is the question of why the \emph{ordering} is what it is.

\paragraph{An OOD-input reading is not enough.} A bare distribution-shift reading---``the classifier was trained against global statistics, anything else is OOD''---would predict that all three substitution conditions hurt similarly, since all three move the activation distribution off the training distribution. The data is incompatible with this prediction in a robust way: same-class is consistently the most damaging condition (Tables~\ref{tab:primary},~\ref{tab:cross_arch}), and on CIFAR-10 it is more damaging than wrong-class by more than $18\times$ in the SmallResNet case. Whatever same-class statistics are doing, they are doing it more thoroughly than a generic OOD perturbation.

\paragraph{An interpretation: BN statistics as an implicit calibrator.} One interpretation consistent with the data is that the trained classification head depends on a tight alignment between the running statistics and its learned weights, and that same-class statistics shift the activation distribution in the direction that the head is most sensitive to for that class---because those statistics encode the class-mean activation pattern. Wrong-class and random-class statistics shift it in directions less specifically aligned with the head's class-$c$ template. Linear refitting (the probe) absorbs whatever offset the substituted statistics induce, which is why probe accuracy is high in every condition. In this reading, BN running statistics function as an output \emph{calibrator} that couples the encoder to the head, and same-class statistics maximally decouple them. We emphasize that this is one possible interpretation rather than a proven mechanism; alternatives---e.g., that same-class normalization specifically collapses the between-class signal that the head exploits---are also consistent with the observations and may be complementary.

\paragraph{The non-monotone probe.} Under any of these interpretations, the non-monotone probe (Section~\ref{sec:interp}) is informative: it indicates that the same-class endpoint is not just ``OOD'' relative to global, but a qualitatively different organization of the activation space (one in which within-class variation is suppressed and the features are very linearly separable). The intermediate regime mixes the two and is worse for linear separability than either endpoint.

\paragraph{Connection to prior empirical work.} The interpretation also gives a perspective on why methods like AdaBN \citep{li2018revisiting} and TENT \citep{wang2021tent} work: they re-estimate BN statistics on data the classifier head must work with, recalibrating the encoder-to-head coupling rather than learning new features. Our results predict that their benefit is largely orthogonal to representation quality, and that pairing them with classifier-head adjustment should magnify their effect. The same lens applies to several other practical settings: BN-folding at inference fuses the running statistics into the convolution, making the calibrator-head coupling an architectural property rather than a runtime one; normalization-free designs such as NFNets~\citep{brock2021highperformance} avoid the coupling entirely; and post-hoc model compression that quantizes or prunes BN buffers should be evaluated for the kind of asymmetric sensitivity we document, since seemingly small statistic perturbations can have outsized effects on accuracy without affecting the underlying representation.

\paragraph{Methodological implication.} A network at 0.17\% accuracy on Tiny-ImageNet-200 may have a 100\%-linearly-separable representation. Reading classification accuracy as a proxy for representation quality is unsafe in any setting in which the encoder-to-head alignment can be perturbed independently of the encoder---domain shift, adversarial inputs, normalization-layer interventions. Diagnosing such failures benefits from separately measuring representation quality (e.g., via a linear probe) and head-alignment (e.g., via ECE plus a temperature-scaling check).

\section{Limitations and Conclusion}
\label{sec:limitations}

\paragraph{Limitations.} (i) All experiments are on natural-image classification at scales up to Tiny-ImageNet-200. Whether the asymmetry holds on ImageNet-1k, fine-grained recognition, or non-vision domains is open. (ii) Architectures with non-BN normalization (LayerNorm, RMSNorm, transformers) carry no running statistics and are not directly addressed. (iii) Wrong-class is operationalized via the fixed mapping $c \mapsto (c+\lfloor C/2\rfloor) \bmod C$; characterizing the full $C\times C$ surface of stat-class pairings would be a richer analysis. (iv) We replace all BN layers simultaneously; layer-wise replacement could reveal which depths are most sensitive. (v) When we replace running statistics we leave the affine parameters $(\gamma_\ell,\beta_\ell)$ unchanged; cleanly disentangling their contribution is left to future work. (vi) ECE with fixed binning has known sensitivities \citep{nixon2019measuring}; we use 15 bins following standard practice.

\paragraph{Conclusion.} Substituting class-conditional BN running statistics at inference produces a sharp asymmetric ordering---same $\ll$ random/wrong $\ll$ global---that holds across four architectures and three datasets. The encoder remains class-discriminative under the same-class substitution (linear probe $\geq 99.7\%$), and the failure is neither a generic normalization-swap effect (GroupNorm models are unaffected) nor a logit-scale issue (temperature scaling halves ECE but recovers 0\% accuracy). A linear interpolation between class-conditional and global statistics shows a smooth monotone recovery of classifier accuracy and a non-monotone probe minimum near $\alpha\approx 0.6$. We treat this as a forensic observation about a specific intervention rather than a sweeping claim about the role of BN statistics, and we discuss interpretations consistent with it. The data we report constrain---rather than determine---a mechanistic story; we hope the asymmetry itself becomes a useful target for theoretical work on BN.

\paragraph{Reproducibility.} Code and data are available at \url{https://anonymous.4open.science/r/bn-calibration}.

\bibliographystyle{abbrvnat}
\bibliography{refs}

\newpage
\appendix

\section{Experimental configuration}
\label{app:config}

\begin{table}[h]
\caption{Full experimental configuration.}
\label{tab:config}
\centering
\small
\begin{tabular}{@{}ll@{}}
\toprule
Parameter & Value \\
\midrule
\textbf{SmallResNet} & BasicBlock, $[2,2,2,2]$, channels $[32,64,128,256]$, 20 BN layers \\
\textbf{SmallResNet64} & BasicBlock, 5-stage, channels $[32,64,128,256,512]$, for $64{\times}64$ \\
\textbf{VGG-11-BN} & Standard VGG-11 with BN, 8 BN layers \\
\textbf{SimpleCNN} & 4-layer CNN, 4 BN layers \\
\textbf{WideResNet-28-10} & depth=28, widen\_factor=10, dropout=0.3 \\
Training epochs & 20 (CIFAR-10), 25 (CIFAR-100, Tiny-ImageNet-200) \\
Batch size & 128 \\
Optimizer & SGD (momentum=0.9, weight\_decay=$5{\times}10^{-4}$) \\
Learning rate & 0.1, cosine annealing \\
Augmentation & RandomCrop(32, pad=4), RandomHorizontalFlip (CIFAR) \\
Normalization (input) & $\mu{=}(0.4914, 0.4822, 0.4465)$, $\sigma{=}(0.2470, 0.2435, 0.2616)$ \\
Linear probe & Ridge ($\lambda{=}10^{-3}$, CIFAR-10) or LogisticRegression ($\geq 100$-class), 80/20 split \\
ECE & 15 equal-width bins \\
Temperature scaling & NLL-optimized scalar $T$ on validation set \\
Seeds & 42--51 (SR/CIFAR-10), 42--46 (SR/CIFAR-100), 42--44 (others) \\
Device & Apple Silicon (MPS backend) \\
\bottomrule
\end{tabular}
\end{table}

\section{Per-class accuracy deltas}
\label{app:perclass}
On CIFAR-10/SmallResNet, per-class same-minus-global accuracy deltas (averaged over 10 seeds) range from $-0.97$ (\emph{automobile}) to $-0.72$ (\emph{cat}). All ten classes show deltas $\leq -0.72$, ruling out the hypothesis that the collapse is concentrated in a small subset of classes.

\section{Cross-architecture results in detail}
\label{app:cross_arch}

\begin{table}[h]
\caption{Cross-architecture results on CIFAR-10 with all four conditions and the linear probe (mean $\pm$ std, $n{=}3$ seeds for VGG/SimpleCNN/WRN).}
\label{tab:cross_arch_full}
\centering
\small
\begin{tabular}{@{}llcccc@{}}
\toprule
Architecture & Condition & Acc.\ (\%) & Conf.\ (\%) & Probe (\%) \\
\midrule
\multirow{4}{*}{VGG-11-BN}
  & Global       & $85.44 \pm 1.28$ & $80.47 \pm 1.78$ & $85.00 \pm 1.80$ \\
  & Same-class   & $10.88 \pm 1.91$ & $10.29 \pm 1.23$ & $98.37 \pm 0.63$ \\
  & Wrong-class  & $62.96 \pm 1.25$ & $58.61 \pm 1.31$ & $99.57 \pm 0.27$ \\
  & Random-class & $57.50 \pm 3.03$ & $52.29 \pm 3.12$ & $95.28 \pm 0.78$ \\
\midrule
\multirow{4}{*}{SimpleCNN}
  & Global       & $83.17 \pm 0.23$ & $75.68 \pm 0.11$ & $80.53 \pm 1.44$ \\
  & Same-class   & $26.33 \pm 0.48$ & $20.71 \pm 0.18$ & $54.68 \pm 0.74$ \\
  & Wrong-class  & $78.10 \pm 0.46$ & $65.79 \pm 0.59$ & $93.67 \pm 0.27$ \\
  & Random-class & $72.11 \pm 0.36$ & $59.20 \pm 0.33$ & $90.82 \pm 0.52$ \\
\midrule
\multirow{4}{*}{WRN-28-10}
  & Global       & $90.61 \pm 0.51$ & $89.36 \pm 0.43$ & $91.02 \pm 0.28$ \\
  & Same-class   & $11.12 \pm 0.90$ & $11.04 \pm 0.73$ & $99.78 \pm 0.09$ \\
  & Wrong-class  & $64.06 \pm 3.14$ & $52.64 \pm 3.51$ & $99.60 \pm 0.07$ \\
  & Random-class & $62.41 \pm 1.76$ & $51.21 \pm 2.25$ & $96.63 \pm 0.44$ \\
\bottomrule
\end{tabular}
\end{table}

\section{Cross-dataset results in detail}
\label{app:cross_data}

\begin{table}[h]
\caption{Full four-condition results on CIFAR-100 (SmallResNet, $n{=}5$ seeds, logistic-regression probe).}
\label{tab:cifar100}
\centering
\small
\begin{tabular}{@{}lccc@{}}
\toprule
Condition & Acc.\ (\%) & Conf.\ (\%) & Probe (\%) \\
\midrule
Global       & $70.11 \pm 0.57$  & $63.98 \pm 0.73$  & $64.23 \pm 1.34$   \\
Same-class   & $0.60 \pm 0.49$   & $0.94 \pm 0.00$   & $100.00 \pm 0.00$  \\
Wrong-class  & $1.00 \pm 0.00$   & $1.00 \pm 0.00$   & $100.00 \pm 0.00$  \\
Random-class & $1.00 \pm 0.63$   & $1.00 \pm 0.00$   & $65.49 \pm 8.75$   \\
\bottomrule
\end{tabular}
\end{table}

\begin{table}[h]
\caption{Full four-condition results on Tiny-ImageNet-200 (SmallResNet64, $n{=}3$ seeds, logistic-regression probe).}
\label{tab:tinyimagenet}
\centering
\small
\begin{tabular}{@{}lccc@{}}
\toprule
Condition & Acc.\ (\%) & Conf.\ (\%) & Probe (\%) \\
\midrule
Global       & $60.77 \pm 0.30$  & $53.75 \pm 0.24$  & $51.83 \pm 0.55$   \\
Same-class   & $0.17 \pm 0.24$   & $0.49 \pm 0.00$   & $100.00 \pm 0.00$  \\
Wrong-class  & $0.50 \pm 0.00$   & $0.50 \pm 0.00$   & $100.00 \pm 0.00$  \\
Random-class & $0.50 \pm 0.00$   & $0.50 \pm 0.00$   & $57.02 \pm 3.41$   \\
\bottomrule
\end{tabular}
\end{table}

\section{Full interpolation sweep}
\label{app:interp}

\begin{table}[h]
\caption{11-point interpolation from same-class ($\alpha=0$) to global ($\alpha=1$) statistics on SmallResNet/CIFAR-10.}
\label{tab:interp_full}
\centering
\small
\begin{tabular}{@{}ccc@{}}
\toprule
$\alpha$ & Accuracy (\%) & Linear Probe (\%) \\
\midrule
0.0 & 3.44  & 99.90 \\
0.1 & 7.13  & 99.70 \\
0.2 & 13.41 & 98.67 \\
0.3 & 23.14 & 94.85 \\
0.4 & 36.47 & 87.10 \\
0.5 & 51.61 & 77.60 \\
0.6 & 65.17 & 74.65 \\
0.7 & 75.26 & 77.88 \\
0.8 & 82.62 & 83.00 \\
0.9 & 87.64 & 87.53 \\
1.0 & 91.01 & 90.85 \\
\bottomrule
\end{tabular}
\end{table}

\section{Detailed ECE and temperature-scaling sweep}
\label{app:ece}

\begin{table}[h]
\caption{Expected Calibration Error by condition (SmallResNet, CIFAR-10, $n{=}3$ seeds).}
\label{tab:ece_full}
\centering
\small
\begin{tabular}{@{}lcc@{}}
\toprule
Condition & Accuracy (\%) & ECE \\
\midrule
Global (baseline) & $91.33 \pm 0.12$ & $0.025 \pm 0.001$ \\
Same-class        & $2.87 \pm 0.49$  & $0.394 \pm 0.013$ \\
Wrong-class       & $63.40 \pm 4.26$ & $0.064 \pm 0.013$ \\
\bottomrule
\end{tabular}
\end{table}

\begin{table}[h]
\caption{Per-seed temperature-scaling on the same-class condition.}
\label{tab:tempscale_full}
\centering
\small
\begin{tabular}{@{}ccccc@{}}
\toprule
Seed & Glob.\ ECE & SC ECE & SC ECE (post-$T$) & $T$ \\
\midrule
42 & 0.026 & 0.382 & 0.150 & 3.456 \\
43 & 0.024 & 0.412 & 0.148 & 3.428 \\
44 & 0.026 & 0.388 & 0.154 & 3.416 \\
\bottomrule
\end{tabular}
\end{table}

Note that wrong-class statistics produce moderate ECE (0.064) despite substantially reduced accuracy: the network's confidence and accuracy remain roughly aligned, just at a lower level. Same-class statistics produce severely miscalibrated outputs (ECE 0.394) and accuracy near floor. The optimal temperature $T\approx 3.43$ indicates the same-class logits are substantially over-scaled.

\section{Per-seed tables}
\label{app:perseed}

\begin{table}[h]
\caption{SmallResNet CIFAR-10 -- Classification accuracy by seed.}
\label{tab:perseed_acc}
\centering
\small
\begin{tabular}{@{}ccccc@{}}
\toprule
Seed & Global & Same-Class & Wrong-Class & Random-Class \\
\midrule
42 & 90.96 & 3.85 & 60.19 & 55.26 \\
43 & 91.05 & 4.02 & 70.70 & 67.73 \\
44 & 91.01 & 2.46 & 69.93 & 55.45 \\
45 & 91.07 & 4.06 & 64.59 & 56.74 \\
46 & 91.26 & 3.84 & 57.08 & 56.07 \\
47 & 91.30 & 3.37 & 60.49 & 57.01 \\
48 & 90.95 & 2.39 & 67.19 & 50.82 \\
49 & 91.19 & 2.90 & 63.57 & 56.01 \\
50 & 90.93 & 4.22 & 69.00 & 66.09 \\
51 & 91.04 & 4.35 & 70.98 & 68.72 \\
\bottomrule
\end{tabular}
\end{table}

\begin{table}[h]
\caption{SmallResNet CIFAR-10 -- Linear probe accuracy by seed.}
\label{tab:perseed_probe}
\centering
\small
\begin{tabular}{@{}ccccc@{}}
\toprule
Seed & Global & Same-Class & Wrong-Class & Random-Class \\
\midrule
42 & 90.30 & 99.90  & 100.00 & 96.30 \\
43 & 90.65 & 99.80  & 100.00 & 96.85 \\
44 & 91.10 & 99.90  & 100.00 & 97.05 \\
45 & 89.65 & 99.90  & 100.00 & 96.90 \\
46 & 90.80 & 100.00 & 100.00 & 98.30 \\
47 & 91.80 & 99.95  & 100.00 & 99.35 \\
48 & 90.80 & 99.90  & 100.00 & 96.80 \\
49 & 91.15 & 99.90  & 100.00 & 98.90 \\
50 & 91.00 & 99.65  & 100.00 & 96.65 \\
51 & 91.30 & 99.90  & 100.00 & 99.15 \\
\bottomrule
\end{tabular}
\end{table}

\begin{table}[h]
\caption{SmallResNet CIFAR-100 -- Per-seed results ($n{=}5$).}
\label{tab:perseed_cifar100}
\centering
\small
\begin{tabular}{@{}ccccc@{}}
\toprule
Seed & Glob.\ Acc & SC Acc & Glob.\ Probe & SC Probe \\
\midrule
42 & 70.95 & 1.00 & 65.65 & 100.00 \\
43 & 69.68 & 1.00 & 64.85 & 100.00 \\
44 & 70.21 & 0.00 & 65.00 & 100.00 \\
45 & 69.31 & 1.00 & 61.80 & 100.00 \\
46 & 70.42 & 0.00 & 63.85 & 100.00 \\
\bottomrule
\end{tabular}
\end{table}

\begin{table}[h]
\caption{WideResNet-28-10 CIFAR-10 -- Per-seed accuracy ($n{=}3$).}
\label{tab:perseed_wrn_acc}
\centering
\small
\begin{tabular}{@{}ccccc@{}}
\toprule
Seed & Global & Same-Class & Wrong-Class & Random-Class \\
\midrule
42 & 90.17 & 10.44 & 64.21 & 62.49 \\
43 & 90.35 & 12.40 & 67.83 & 64.53 \\
44 & 91.32 & 10.53 & 60.15 & 60.22 \\
\bottomrule
\end{tabular}
\end{table}

\begin{table}[h]
\caption{WideResNet-28-10 CIFAR-10 -- Per-seed linear probe ($n{=}3$).}
\label{tab:perseed_wrn_probe}
\centering
\small
\begin{tabular}{@{}ccccc@{}}
\toprule
Seed & Global & Same-Class & Wrong-Class & Random-Class \\
\midrule
42 & 91.40 & 99.85 & 99.70 & 96.30 \\
43 & 90.90 & 99.85 & 99.55 & 96.35 \\
44 & 90.75 & 99.65 & 99.55 & 97.25 \\
\bottomrule
\end{tabular}
\end{table}

\begin{table}[h]
\caption{SmallResNet64 Tiny-ImageNet-200 -- Per-seed results ($n{=}3$).}
\label{tab:perseed_tinyim}
\centering
\small
\begin{tabular}{@{}cccccc@{}}
\toprule
Seed & Glob.\ Acc & Glob.\ Conf & Glob.\ Probe & SC Acc & SC Probe \\
\midrule
42 & 60.88 & 53.78 & 51.55 & 0.00 & 100.00 \\
43 & 61.08 & 54.04 & 52.60 & 0.50 & 100.00 \\
44 & 60.36 & 53.44 & 51.35 & 0.00 & 100.00 \\
\bottomrule
\end{tabular}
\end{table}

\newpage

\section*{NeurIPS Paper Checklist}

\begin{enumerate}

\item {\bf Claims}
    \item[] Question: Do the main claims made in the abstract and introduction accurately reflect the paper's contributions and scope?
    \item[] Answer: \answerYes{}
    \item[] Justification: The abstract and Section~\ref{sec:intro} state the asymmetric ordering, the architectures and datasets covered, and the linear probe / GroupNorm / temperature-scaling controls; Section~\ref{sec:results} reports the corresponding numbers.
    \item[] Guidelines:
    \begin{itemize}
        \item The answer \answerNA{} means that the abstract and introduction do not include the claims made in the paper.
        \item The abstract and/or introduction should clearly state the claims made, including the contributions made in the paper and important assumptions and limitations. A \answerNo{} or \answerNA{} answer to this question will not be perceived well by the reviewers.
        \item The claims made should match theoretical and experimental results, and reflect how much the results can be expected to generalize to other settings.
        \item It is fine to include aspirational goals as motivation as long as it is clear that these goals are not attained by the paper.
    \end{itemize}

\item {\bf Limitations}
    \item[] Question: Does the paper discuss the limitations of the work performed by the authors?
    \item[] Answer: \answerYes{}
    \item[] Justification: Section~\ref{sec:limitations} enumerates six limitations including dataset scale (up to Tiny-ImageNet-200), the BN-only scope, the fixed wrong-class mapping, the all-layers-at-once substitution, the unchanged affine parameters, and ECE binning sensitivity.
    \item[] Guidelines:
    \begin{itemize}
        \item The answer \answerNA{} means that the paper has no limitation while the answer \answerNo{} means that the paper has limitations, but those are not discussed in the paper.
        \item The authors are encouraged to create a separate ``Limitations'' section in their paper.
        \item The paper should point out any strong assumptions and how robust the results are to violations of these assumptions (e.g., independence assumptions, noiseless settings, model well-specification, asymptotic approximations only holding locally). The authors should reflect on how these assumptions might be violated in practice and what the implications would be.
        \item The authors should reflect on the scope of the claims made, e.g., if the approach was only tested on a few datasets or with a few runs. In general, empirical results often depend on implicit assumptions, which should be articulated.
        \item The authors should reflect on the factors that influence the performance of the approach. For example, a facial recognition algorithm may perform poorly when image resolution is low or images are taken in low lighting. Or a speech-to-text system might not be used reliably to provide closed captions for online lectures because it fails to handle technical jargon.
        \item The authors should discuss the computational efficiency of the proposed algorithms and how they scale with dataset size.
        \item If applicable, the authors should discuss possible limitations of their approach to address problems of privacy and fairness.
        \item While the authors might fear that complete honesty about limitations might be used by reviewers as grounds for rejection, a worse outcome might be that reviewers discover limitations that aren't acknowledged in the paper. The authors should use their best judgment and recognize that individual actions in favor of transparency play an important role in developing norms that preserve the integrity of the community. Reviewers will be specifically instructed to not penalize honesty concerning limitations.
    \end{itemize}

\item {\bf Theory assumptions and proofs}
    \item[] Question: For each theoretical result, does the paper provide the full set of assumptions and a complete (and correct) proof?
    \item[] Answer: \answerNA{}
    \item[] Justification: The paper is an empirical forensic study and contains no theorems or formal proofs.
    \item[] Guidelines:
    \begin{itemize}
        \item The answer \answerNA{} means that the paper does not include theoretical results.
        \item All the theorems, formulas, and proofs in the paper should be numbered and cross-referenced.
        \item All assumptions should be clearly stated or referenced in the statement of any theorems.
        \item The proofs can either appear in the main paper or the supplemental material, but if they appear in the supplemental material, the authors are encouraged to provide a short proof sketch to provide intuition.
        \item Inversely, any informal proof provided in the core of the paper should be complemented by formal proofs provided in appendix or supplemental material.
        \item Theorems and Lemmas that the proof relies upon should be properly referenced.
    \end{itemize}

    \item {\bf Experimental result reproducibility}
    \item[] Question: Does the paper fully disclose all the information needed to reproduce the main experimental results of the paper to the extent that it affects the main claims and/or conclusions of the paper (regardless of whether the code and data are provided or not)?
    \item[] Answer: \answerYes{}
    \item[] Justification: Section~\ref{sec:method} and Appendix~\ref{app:config} fully specify architectures, optimizer, learning-rate schedule, augmentation, batch size, epochs, seeds (42--51 for the headline configuration), the substitution procedure, probe protocol, and ECE/temperature-scaling protocol; anonymized code is released at the URL in Section~\ref{sec:limitations}.
    \item[] Guidelines:
    \begin{itemize}
        \item The answer \answerNA{} means that the paper does not include experiments.
        \item If the paper includes experiments, a \answerNo{} answer to this question will not be perceived well by the reviewers: Making the paper reproducible is important, regardless of whether the code and data are provided or not.
        \item If the contribution is a dataset and\slash or model, the authors should describe the steps taken to make their results reproducible or verifiable.
        \item Depending on the contribution, reproducibility can be accomplished in various ways. For example, if the contribution is a novel architecture, describing the architecture fully might suffice, or if the contribution is a specific model and empirical evaluation, it may be necessary to either make it possible for others to replicate the model with the same dataset, or provide access to the model. In general. releasing code and data is often one good way to accomplish this, but reproducibility can also be provided via detailed instructions for how to replicate the results, access to a hosted model (e.g., in the case of a large language model), releasing of a model checkpoint, or other means that are appropriate to the research performed.
        \item While NeurIPS does not require releasing code, the conference does require all submissions to provide some reasonable avenue for reproducibility, which may depend on the nature of the contribution. For example
        \begin{enumerate}
            \item If the contribution is primarily a new algorithm, the paper should make it clear how to reproduce that algorithm.
            \item If the contribution is primarily a new model architecture, the paper should describe the architecture clearly and fully.
            \item If the contribution is a new model (e.g., a large language model), then there should either be a way to access this model for reproducing the results or a way to reproduce the model (e.g., with an open-source dataset or instructions for how to construct the dataset).
            \item We recognize that reproducibility may be tricky in some cases, in which case authors are welcome to describe the particular way they provide for reproducibility. In the case of closed-source models, it may be that access to the model is limited in some way (e.g., to registered users), but it should be possible for other researchers to have some path to reproducing or verifying the results.
        \end{enumerate}
    \end{itemize}


\item {\bf Open access to data and code}
    \item[] Question: Does the paper provide open access to the data and code, with sufficient instructions to faithfully reproduce the main experimental results, as described in supplemental material?
    \item[] Answer: \answerYes{}
    \item[] Justification: Anonymized code is available at \url{https://anonymous.4open.science/r/bn-calibration} as stated in Section~\ref{sec:limitations}; all datasets used (CIFAR-10, CIFAR-100, Tiny-ImageNet-200) are public.
    \item[] Guidelines:
    \begin{itemize}
        \item The answer \answerNA{} means that paper does not include experiments requiring code.
        \item Please see the NeurIPS code and data submission guidelines (\url{https://neurips.cc/public/guides/CodeSubmissionPolicy}) for more details.
        \item While we encourage the release of code and data, we understand that this might not be possible, so \answerNo{} is an acceptable answer. Papers cannot be rejected simply for not including code, unless this is central to the contribution (e.g., for a new open-source benchmark).
        \item The instructions should contain the exact command and environment needed to run to reproduce the results. See the NeurIPS code and data submission guidelines (\url{https://neurips.cc/public/guides/CodeSubmissionPolicy}) for more details.
        \item The authors should provide instructions on data access and preparation, including how to access the raw data, preprocessed data, intermediate data, and generated data, etc.
        \item The authors should provide scripts to reproduce all experimental results for the new proposed method and baselines. If only a subset of experiments are reproducible, they should state which ones are omitted from the script and why.
        \item At submission time, to preserve anonymity, the authors should release anonymized versions (if applicable).
        \item Providing as much information as possible in supplemental material (appended to the paper) is recommended, but including URLs to data and code is permitted.
    \end{itemize}


\item {\bf Experimental setting/details}
    \item[] Question: Does the paper specify all the training and test details (e.g., data splits, hyperparameters, how they were chosen, type of optimizer) necessary to understand the results?
    \item[] Answer: \answerYes{}
    \item[] Justification: Section~\ref{sec:method} specifies architectures, optimizer (SGD, momentum 0.9, weight decay $5\times 10^{-4}$), learning-rate schedule (cosine), augmentation, epochs, and the four inference conditions; Appendix~\ref{app:config} consolidates the full configuration.
    \item[] Guidelines:
    \begin{itemize}
        \item The answer \answerNA{} means that the paper does not include experiments.
        \item The experimental setting should be presented in the core of the paper to a level of detail that is necessary to appreciate the results and make sense of them.
        \item The full details can be provided either with the code, in appendix, or as supplemental material.
    \end{itemize}

\item {\bf Experiment statistical significance}
    \item[] Question: Does the paper report error bars suitably and correctly defined or other appropriate information about the statistical significance of the experiments?
    \item[] Answer: \answerYes{}
    \item[] Justification: All result tables report mean $\pm$ std over multiple seeds (10 for the headline SmallResNet/CIFAR-10 configuration, 5 for SmallResNet/CIFAR-100, 3 for the others); paired $t$-tests are reported in Section~\ref{sec:ordering} ($t(9)=-374.1$, $p=3.5\times 10^{-20}$ for same-vs-global; $t(9)=-39.1$, $p<10^{-10}$ for same-vs-wrong).
    \item[] Guidelines:
    \begin{itemize}
        \item The answer \answerNA{} means that the paper does not include experiments.
        \item The authors should answer \answerYes{} if the results are accompanied by error bars, confidence intervals, or statistical significance tests, at least for the experiments that support the main claims of the paper.
        \item The factors of variability that the error bars are capturing should be clearly stated (for example, train/test split, initialization, random drawing of some parameter, or overall run with given experimental conditions).
        \item The method for calculating the error bars should be explained (closed form formula, call to a library function, bootstrap, etc.)
        \item The assumptions made should be given (e.g., Normally distributed errors).
        \item It should be clear whether the error bar is the standard deviation or the standard error of the mean.
        \item It is OK to report 1-sigma error bars, but one should state it. The authors should preferably report a 2-sigma error bar than state that they have a 96\% CI, if the hypothesis of Normality of errors is not verified.
        \item For asymmetric distributions, the authors should be careful not to show in tables or figures symmetric error bars that would yield results that are out of range (e.g., negative error rates).
        \item If error bars are reported in tables or plots, the authors should explain in the text how they were calculated and reference the corresponding figures or tables in the text.
    \end{itemize}

\item {\bf Experiments compute resources}
    \item[] Question: For each experiment, does the paper provide sufficient information on the computer resources (type of compute workers, memory, time of execution) needed to reproduce the experiments?
    \item[] Answer: \answerYes{}
    \item[] Justification: Appendix~\ref{app:config} lists the device used (Apple Silicon, MPS backend); per-run training is short (20--25 epochs at batch size 128 on CIFAR-scale data), and the substitution / probe / temperature-scaling pipelines are inference-only and inexpensive.
    \item[] Guidelines:
    \begin{itemize}
        \item The answer \answerNA{} means that the paper does not include experiments.
        \item The paper should indicate the type of compute workers CPU or GPU, internal cluster, or cloud provider, including relevant memory and storage.
        \item The paper should provide the amount of compute required for each of the individual experimental runs as well as estimate the total compute.
        \item The paper should disclose whether the full research project required more compute than the experiments reported in the paper (e.g., preliminary or failed experiments that didn't make it into the paper).
    \end{itemize}

\item {\bf Code of ethics}
    \item[] Question: Does the research conducted in the paper conform, in every respect, with the NeurIPS Code of Ethics \url{https://neurips.cc/public/EthicsGuidelines}?
    \item[] Answer: \answerYes{}
    \item[] Justification: The work uses only public benchmark image-classification datasets, releases no high-risk models or scraped data, and preserves anonymity in the submission.
    \item[] Guidelines:
    \begin{itemize}
        \item The answer \answerNA{} means that the authors have not reviewed the NeurIPS Code of Ethics.
        \item If the authors answer \answerNo, they should explain the special circumstances that require a deviation from the Code of Ethics.
        \item The authors should make sure to preserve anonymity (e.g., if there is a special consideration due to laws or regulations in their jurisdiction).
    \end{itemize}


\item {\bf Broader impacts}
    \item[] Question: Does the paper discuss both potential positive societal impacts and negative societal impacts of the work performed?
    \item[] Answer: \answerNA{}
    \item[] Justification: The work is a forensic empirical study of an inference-time intervention on existing image classifiers; it does not propose a deployable system or release new capabilities, so direct societal impact is minimal.
    \item[] Guidelines:
    \begin{itemize}
        \item The answer \answerNA{} means that there is no societal impact of the work performed.
        \item If the authors answer \answerNA{} or \answerNo, they should explain why their work has no societal impact or why the paper does not address societal impact.
        \item Examples of negative societal impacts include potential malicious or unintended uses (e.g., disinformation, generating fake profiles, surveillance), fairness considerations (e.g., deployment of technologies that could make decisions that unfairly impact specific groups), privacy considerations, and security considerations.
        \item The conference expects that many papers will be foundational research and not tied to particular applications, let alone deployments. However, if there is a direct path to any negative applications, the authors should point it out. For example, it is legitimate to point out that an improvement in the quality of generative models could be used to generate Deepfakes for disinformation. On the other hand, it is not needed to point out that a generic algorithm for optimizing neural networks could enable people to train models that generate Deepfakes faster.
        \item The authors should consider possible harms that could arise when the technology is being used as intended and functioning correctly, harms that could arise when the technology is being used as intended but gives incorrect results, and harms following from (intentional or unintentional) misuse of the technology.
        \item If there are negative societal impacts, the authors could also discuss possible mitigation strategies (e.g., gated release of models, providing defenses in addition to attacks, mechanisms for monitoring misuse, mechanisms to monitor how a system learns from feedback over time, improving the efficiency and accessibility of ML).
    \end{itemize}

\item {\bf Safeguards}
    \item[] Question: Does the paper describe safeguards that have been put in place for responsible release of data or models that have a high risk for misuse (e.g., pre-trained language models, image generators, or scraped datasets)?
    \item[] Answer: \answerNA{}
    \item[] Justification: No high-risk models, generative systems, or scraped datasets are released; only analysis code and small image-classification checkpoints trained on standard public benchmarks.
    \item[] Guidelines:
    \begin{itemize}
        \item The answer \answerNA{} means that the paper poses no such risks.
        \item Released models that have a high risk for misuse or dual-use should be released with necessary safeguards to allow for controlled use of the model, for example by requiring that users adhere to usage guidelines or restrictions to access the model or implementing safety filters.
        \item Datasets that have been scraped from the Internet could pose safety risks. The authors should describe how they avoided releasing unsafe images.
        \item We recognize that providing effective safeguards is challenging, and many papers do not require this, but we encourage authors to take this into account and make a best faith effort.
    \end{itemize}

\item {\bf Licenses for existing assets}
    \item[] Question: Are the creators or original owners of assets (e.g., code, data, models), used in the paper, properly credited and are the license and terms of use explicitly mentioned and properly respected?
    \item[] Answer: \answerYes{}
    \item[] Justification: CIFAR-10 and CIFAR-100 \citep{ioffe2015batch} (cited via standard sources) and Tiny-ImageNet-200 are public benchmark datasets used under their standard research-use terms; original BN, GroupNorm, WideResNet, AdaBN, TENT, and temperature-scaling references are cited in Section~\ref{sec:bg}.
    \item[] Guidelines:
    \begin{itemize}
        \item The answer \answerNA{} means that the paper does not use existing assets.
        \item The authors should cite the original paper that produced the code package or dataset.
        \item The authors should state which version of the asset is used and, if possible, include a URL.
        \item The name of the license (e.g., CC-BY 4.0) should be included for each asset.
        \item For scraped data from a particular source (e.g., website), the copyright and terms of service of that source should be provided.
        \item If assets are released, the license, copyright information, and terms of use in the package should be provided. For popular datasets, \url{paperswithcode.com/datasets} has curated licenses for some datasets. Their licensing guide can help determine the license of a dataset.
        \item For existing datasets that are re-packaged, both the original license and the license of the derived asset (if it has changed) should be provided.
        \item If this information is not available online, the authors are encouraged to reach out to the asset's creators.
    \end{itemize}

\item {\bf New assets}
    \item[] Question: Are new assets introduced in the paper well documented and is the documentation provided alongside the assets?
    \item[] Answer: \answerYes{}
    \item[] Justification: We release the analysis and training code anonymously at the URL given in Section~\ref{sec:limitations}; the repository contains a README documenting how to reproduce each table.
    \item[] Guidelines:
    \begin{itemize}
        \item The answer \answerNA{} means that the paper does not release new assets.
        \item Researchers should communicate the details of the dataset\slash code\slash model as part of their submissions via structured templates. This includes details about training, license, limitations, etc.
        \item The paper should discuss whether and how consent was obtained from people whose asset is used.
        \item At submission time, remember to anonymize your assets (if applicable). You can either create an anonymized URL or include an anonymized zip file.
    \end{itemize}

\item {\bf Crowdsourcing and research with human subjects}
    \item[] Question: For crowdsourcing experiments and research with human subjects, does the paper include the full text of instructions given to participants and screenshots, if applicable, as well as details about compensation (if any)?
    \item[] Answer: \answerNA{}
    \item[] Justification: The paper does not involve crowdsourcing or research with human subjects.
    \item[] Guidelines:
    \begin{itemize}
        \item The answer \answerNA{} means that the paper does not involve crowdsourcing nor research with human subjects.
        \item Including this information in the supplemental material is fine, but if the main contribution of the paper involves human subjects, then as much detail as possible should be included in the main paper.
        \item According to the NeurIPS Code of Ethics, workers involved in data collection, curation, or other labor should be paid at least the minimum wage in the country of the data collector.
    \end{itemize}

\item {\bf Institutional review board (IRB) approvals or equivalent for research with human subjects}
    \item[] Question: Does the paper describe potential risks incurred by study participants, whether such risks were disclosed to the subjects, and whether Institutional Review Board (IRB) approvals (or an equivalent approval/review based on the requirements of your country or institution) were obtained?
    \item[] Answer: \answerNA{}
    \item[] Justification: The paper does not involve human subjects, so IRB approval is not applicable.
    \item[] Guidelines:
    \begin{itemize}
        \item The answer \answerNA{} means that the paper does not involve crowdsourcing nor research with human subjects.
        \item Depending on the country in which research is conducted, IRB approval (or equivalent) may be required for any human subjects research. If you obtained IRB approval, you should clearly state this in the paper.
        \item We recognize that the procedures for this may vary significantly between institutions and locations, and we expect authors to adhere to the NeurIPS Code of Ethics and the guidelines for their institution.
        \item For initial submissions, do not include any information that would break anonymity (if applicable), such as the institution conducting the review.
    \end{itemize}

\item {\bf Declaration of LLM usage}
    \item[] Question: Does the paper describe the usage of LLMs if it is an important, original, or non-standard component of the core methods in this research? Note that if the LLM is used only for writing, editing, or formatting purposes and does \emph{not} impact the core methodology, scientific rigor, or originality of the research, declaration is not required.
    \item[] Answer: \answerNA{}
    \item[] Justification: The core methods (BN-statistic substitution, linear probing, temperature scaling, GroupNorm control) do not involve LLMs in any non-standard or load-bearing role.
    \item[] Guidelines:
    \begin{itemize}
        \item The answer \answerNA{} means that the core method development in this research does not involve LLMs as any important, original, or non-standard components.
        \item Please refer to our LLM policy in the NeurIPS handbook for what should or should not be described.
    \end{itemize}

\end{enumerate}

\end{document}
